\section{Limitations and Future Work}\label{sec:limits}

Several limitations bound the scope of these results. First, the study covers a
\emph{single} reliability zone, NEMA (Northeast Massachusetts / Boston); whether the
same feature design and source-matching discipline transfer to other ISO-NE zones, or
to other markets with different load shapes and weather regimes, remains to be shown.
Second, the live horizon-matched evaluation against ISO-NE spans only about $30$ days
of spring and early summer. This window is short and seasonally narrow: it excludes
peak-summer cooling load and winter heating extremes, the periods where day-ahead
forecasting is most consequential and most difficult. The test-set results
(Section~\ref{sec:results}) draw on a full year of held-out 2025 data and are
correspondingly more robust, but the head-to-head with the operator should be read as
indicative rather than conclusive.

Third, the load series differ slightly between the training source (wholesale-load,
WHLSECOST) and the real-time-demand series used as live actuals; small definitional
differences between these series add noise to the live comparison. Fourth, the
Open-Meteo ERA5 archive does not provide a visibility channel, which was filled with a
constant during training; visibility is therefore inert as a feature. Finally, Beacon
produces \emph{point} forecasts only --- it emits no predictive intervals or quantile
bands, which limits its usefulness for reserve sizing and risk-aware dispatch, where
uncertainty quantification is valuable.

These limitations map onto a clear program of future work. The most immediate
extension is \emph{probabilistic} forecasting: quantile or interval bands around each
per-horizon prediction, so that the forecaster communicates uncertainty rather than a
single number. Beyond point accuracy, ensembling the gradient-boosted models with a
transformer-style sequence model such as Informer~\cite{zhou2021informer} may capture
longer-range temporal structure that the lag features miss. Broader evaluation ---
extending across all ISO-NE zones and across a full seasonal cycle --- would test
generalization directly. Systematic hyper-parameter search via Optuna, and
recency-weighted retraining that adapts to gradual changes in load behavior, are
natural refinements to the training procedure that we leave to future work.
